\documentclass[11pt]{article}
\usepackage[margin=1in]{geometry}
\usepackage{amsmath,amssymb,amsthm,mathtools}
\usepackage[hidelinks]{hyperref}
\setlength{\emergencystretch}{3em}
\allowdisplaybreaks[2]
\newtheorem{proposition}{Proposition}[section]
\newtheorem{lemma}[proposition]{Lemma}
\theoremstyle{remark}\newtheorem{remark}[proposition]{Remark}
\newcommand{\E}{\mathbb E}
\newcommand{\dd}{\,d}
\newcommand{\one}{\mathbf1}
\newcommand{\Gp}{G_+}
\newcommand{\KK}{\mathcal K}
\newcommand{\BB}{\mathcal B}
\title{Strong c-form reduction: selected-gate estimates\\
and the remaining S2 dependency}
\author{Analytical increment for review}
\date{Version 1; October 21, 2026}
\begin{document}
\maketitle
\begin{abstract}
On the stated residual-stopped interval, we bound the primary
cohort's selected cluster-2 regression defect, including its
scaled receiver derivative, by an integrable small-constant
gate error plus a vanishing transverse term. We give exact
non-core moment definitions, a negative-part bound for the
candidate weak forcing, and a radial identity with its remainder.
The requested complete c-form theorem is not closed.
The proposed charge of the secondary cohort invokes a
conditional lemma whose distance-envelope and all-endpoint
gain hypotheses remain open. Neither those hypotheses nor
the resulting secondary energy budget is among the stated
guards. A separate integrated outside-mass cost also remains.
We display these missing inequalities and stop without
assuming them, applying the passage theorem, or starting S2.
\end{abstract}

\section{Scope and the precise stopping point}
Every new label has prefix \texttt{an03cfrv1:}.
Use the original balanced untied Gaussian-mixture flow and the
notation of CLOCK and LCRC:
\[
 P=\tfrac12P_1+\tfrac12P_2,\quad
 P_1=N(e_1,q^2I_2),\quad P_2=N(\rho e_2,q^2I_2),\quad
 \rho\in[13/20,7/10],\quad \lambda=\rho^2,\quad Q=\lambda+q^2,
\]
\[
 U'=R_\theta^TW,\quad W'=R_\theta U,\quad
 R_\theta=\E_P[(X-f)X^T\one_{\{U\cdot X>0\}}],
 \qquad |U|^2=|W|^2=m .
\]
The fixed partition is $\KK,\Gp,\Sigma,\BB$, with
$\mathcal D=\emptyset$. As in SEP,
$\Sigma=\{\beta\in Q2:c_0<Kq^\gamma\}$ with Q2EC's fixed cutoff.
For a cohort $S$, put
\[
 f_S=\int_S W(U\cdot X)_+\dd\lambda_0,\quad
 E_S=\tfrac12\int_S(U_2^2+W_2^2),\quad
 J_S=\tfrac12\int_S(U_1^2+W_1^2),\quad \mu_S=E_S+J_S,
 \quad c_S=Q^{-1}\E_{P_2}[X_2(f_S)_2].
\]
Integrals without an indicated label measure use $\dd\lambda_0$.
Set $G=\Gp\cup\Sigma$, $c=\sum_Sc_S$ and
$I(t)=Q\int_{t_0}^t(1-c)$, with $t_0=t_h$.

The requested interval has $0\le c\le c_b<1$, duration
$T_*=t_*-t_0\le L_T\ell$, $\ell=\log(1/q)$, and guards
\begin{equation}
 J_G<q^2K_Y^2,\qquad \int_{t_0}^{t_*}r<B,\qquad
 E_{\Gp}<E_g,\qquad
 \|e_2\|_{P_1}+\|e_1\|_{P_2}\le M_{\rm res}q .
 \label{an03cfrv1:guards}
\end{equation}
Here $r=\sup_\theta|R_{11}|$, $e=X-f$.
We rename the question's cross-residual constant $M$ to
$M_{\rm res}$ to distinguish it from strong mass.
Loss monotonicity supplies a fixed diagonal residual norm bound.
The now-proved S1 entry contracts supply RCET's initial conditions.
Thus RCET applies to $\Gp$ on this stopped interval, for its
admissible fixed $C_0$, without using a later S2 estimate.

The source dependency that prevents the full conclusion is
\begin{equation}
 \boxed{\quad
 \texttt{LCRC:profile}+\texttt{LCRC:distance}
       +\texttt{LCRC:gain}+\text{stopped clock}
 \ \Longrightarrow\
 E_\Sigma(t)\le C_\Sigma(q^5+q^{\Gamma_\nu}H(t)).
 \quad}
 \label{an03cfrv1:dependency}
\end{equation}
Here $H=H_{\mathcal C_q}$ and
$\Gamma_\nu\ge289169/924000>3/10$.
This is LCRC's \texttt{secondary} lemma, not an implication
of \eqref{an03cfrv1:guards} alone.
SEP supplies \texttt{profile} at $t_0$ and explicitly leaves
\texttt{distance} and \texttt{gain} open. The present task
does not assume those contracts and prohibits proving them.
The estimates below isolate the resulting unpaid terms.
They do not assert that the desired original-flow theorem is false.

\section{A receiver-uniform selected-moment lemma}
Fix $L_c\ge79/51$ before the residual guards.
For $|x|\le L_c$, let
\[
 I_x=\one_{\{a_{qx}\cdot X>0\}},\qquad
 g_S=(f_S)_2-c_SX_2,\qquad
 \mathcal D_S(x)=\E_{P_2}[g_S X_2 I_x].
\]
Write $\|h\|_{1,\infty}=\sup_{|x|\le L_c}(|h|+|h'|)$.
For a scalar independent of the output angle $y$, this controls
its receiver derivative also in the induced $\ell^1$ norm.
Constants $C_{L_c}$ below depend only on $L_c$, the ratio box,
and fixed Gaussian moments.

\begin{lemma}[Any fixed cohort, with no angular or sign assumption]
\label{an03cfrv1:moment}
For every balanced state,
\begin{align}
 |c_S|+\|\mathcal D_S\|_{1,\infty}
 &\le C_{L_c}\big(E_S+q\sqrt{E_SJ_S}\big),\label{an03cfrv1:anycohort}\\
 \|\mathcal D_{\BB}\|_{1,\infty}
 &\le C_{L_c}\mu_{\BB}.\label{an03cfrv1:Bselected}
\end{align}
If all donors in $\KK$ have both angles of absolute value at
most $Rq$, with fixed $R$, then
\begin{equation}
 |c_{\KK}|+\|\mathcal D_{\KK}\|_{1,\infty}
       \le C_{L_c,R}q^2\mu_{\KK}.
 \label{an03cfrv1:Kselected}
\end{equation}
The last assertion is an instantaneous geometric implication,
not a proof of core retention after $t_0$.
\end{lemma}
\begin{proof}
Write $X=(qZ,V)$ under $P_2$, with independent standard normal
$Z$ and $V\sim N(\rho,q^2)$. For
$t_x=\tan(qx)/q$, the receiver gate is $Z+t_xV>0$.
For small $q$, $|t_x|\le2L_c$ and $t_x'=\sec^2(qx)\le2$.
Cauchy--Schwarz in labels gives
\begin{equation}
 |(f_S)_2(X)|
 \le 2\sqrt{E_SJ_S}|X_1|+2E_S|X_2|.
 \label{an03cfrv1:outputbound}
\end{equation}
This proves the unselected and selected value bounds using
bounded Gaussian moments and $Q\ge169/400$.

The boundary formula for
$F_S(x)=\E[(f_S)_2X_2I_x]$ is
\begin{equation}
 F_S'(x)=t_x'\E_V\!\left[
 V^2(f_S)_2(-qt_xV,V)\phi(t_xV)\right].
 \label{an03cfrv1:boundary}
\end{equation}
At the boundary,
\[
 |(f_S)_2(-qt_xV,V)|
 \le2|V|\big(E_S+q|t_x|\sqrt{E_SJ_S}\big).
\]
Bounded third moments of $V$ prove the derivative estimate.
The function $T_2(x)=\E[V^2I_x]$ and its derivative are
bounded by $C_{L_c}$, so subtraction of $c_ST_2$ gives
\eqref{an03cfrv1:anycohort}.
Domination by these polynomial Gaussian bounds justifies
the boundary differentiation; the ReLU output is continuous.

Alternatively $|(f_{\BB})_2(X)|\le\mu_{\BB}|X|$ proves
\eqref{an03cfrv1:Bselected} in the same formulas.
For charted core donors,
$|W_2|\le Rq\sqrt m$, $|U_1|\le\sqrt m$,
$|U_2|\le Rq\sqrt m$.
Their output is bounded by
$C_Rq^2\mu_{\KK}(|Z|+|V|)$ on $P_2$, including on the
receiver boundary. This proves \eqref{an03cfrv1:Kselected}.
It does not discard the core first output: that output
still generates the order-one strong overlap term in MOMC.
\end{proof}

\section{The primary contribution closes under the stated guards}
For $\Gp$ define the exact coefficients and negative-gate output
\[
 A_+=\int_{\Gp}W_2U_2,\qquad B_+=\int_{\Gp}W_2U_1,\qquad
 N_+(X)=\int_{\Gp}W_2(U\cdot X)_-,
 \qquad v_-=(-v)_+ .
\]
RCET gives $W_2>0$, so $N_+\ge0$.
The following are exact identities, not full-gate substitutions:
\begin{equation}
 (f_{\Gp})_2=A_+X_2+B_+X_1+N_+,\quad
 \tau_+=Q^{-1}\E_{P_2}[X_2N_+],\quad c_{\Gp}=A_++\tau_+,
 \quad g_{\Gp}=B_+X_1+N_+-\tau_+X_2 .
 \label{an03cfrv1:linearity}
\end{equation}

\begin{proposition}[Closed primary selected-defect budget]
\label{an03cfrv1:primary}
Under \eqref{an03cfrv1:guards} and the proved primary entry
conditions, take $C_0$ sufficiently small for RCET and,
if necessary, smaller depending on $L_c$ and its fixed
RCET margin constants. There is a nonnegative function
$\mathcal T_+(t)$ such that
\begin{align}
 \|\mathcal D_{\Gp}(t)\|_{1,\infty}
 &\le C_{L_c}\left[q\sqrt{E_{\Gp}J_{\Gp}}
                     +\mu_{\Gp}\mathcal T_+(t)\right],
 \label{an03cfrv1:primarypoint}\\
 \int_{t_0}^{t_*}\mathcal T_+
 &\le C_{\rm RCET}C_0+C_{\rm RCET}q^2T_* .
 \label{an03cfrv1:tailintegral}
\end{align}
In particular
\begin{equation}
 \int_{t_0}^{t_*}\|\mathcal D_{\Gp}\|_{1,\infty}
 \le C_{L_c}C_{\rm RCET}E_g C_0+o_q(1).
 \label{an03cfrv1:primaryintegral}
\end{equation}
The error is uniform over stopped endpoints. Dependence on
$(B,E_g,K_Y,M_{\rm res})$ enters only the coefficient of
$C_0$ and vanishing terms, not $C_{L_c}$.
\end{proposition}
\begin{proof}
Cauchy--Schwarz gives $|B_+|\le2\sqrt{E_{\Gp}J_{\Gp}}$.
The exact selected cross moment is
\[
 \E_{P_2}[X_1X_2I_x]=q\E_V[V\phi(t_xV)].
\]
Its value and derivative are bounded by $C_{L_c}q$.
This controls the first term in \eqref{an03cfrv1:linearity}.

For a donor direction $a_\theta$, Gaussian tails and
Cauchy--Schwarz give
\[
 \E_{P_2}[|X|^2\one_{\{a_\theta\cdot X<0\}}]
 \le C\exp[-c(a_{\theta,2}/q)^2]\quad(a_{\theta,2}>0).
\]
The fourth Gaussian moment is uniformly bounded, and the
half-space probability is at most
$\exp[-\rho^2a_{\theta,2}^2/(2q^2)]$; decreasing $c$ proves
the display. Since $|W_2||U|\le m$, it bounds both the
selected value of $N_+X_2$ and $|\tau_+|$ by the mass
times this tail.

For the derivative there is an additional useful margin.
RCET gives $U_2\ge2z/3$ and
$q|U_1|/U_2\le C(qj/z)$, with
$qj/z\le C(C_0e^{-b_*(t-t_0)/2}+q^2)$.
Choose $C_0$ and then $q$ so that
$q|t_xU_1|\le U_2/2$ for every receiver in the fixed chart.
At its boundary,
\[
 (U\cdot X)_-=[-(U_2-qt_xU_1)V]_+ .
\]
It vanishes for $V\ge0$. For $V<0$ it is bounded by
$C_{L_c}|U||V|$.
Formula \eqref{an03cfrv1:boundary}, now applied to $N_+$,
therefore bounds the derivative by
$C_{L_c}\mu_{\Gp}\E[|V|^3\one_{\{V<0\}}]
\le C_{L_c}\mu_{\Gp}e^{-c/q^2}$.
The subtraction $-\tau_+T_2$ has the already established
value/derivative bound.

RCET's growing margin allows, after enlarging fixed constants,
\[
 \mathcal T_+(t)=C\exp\left[
 -\frac{c}{q^2+C_*^2C_0^2e^{-b_*(t-t_0)}}\right]+Ce^{-c/q^2},
 \qquad b_*=Q(1-c_b)/2.
\]
Splitting the denominator as in RCET, and using
$e^{-c/w^2}\le C_c w$ and $e^{-c/q^2}=O(q^2)$, proves
\eqref{an03cfrv1:tailintegral}.
Finally $\mu_{\Gp}\le E_g+q^2K_Y^2$ and
$q\sqrt{E_{\Gp}J_{\Gp}}\le q^2K_Y\sqrt{E_g}$.
Their integrals prove \eqref{an03cfrv1:primaryintegral}.
No S2 or terminal amplitude cap was used.
\end{proof}

For the total regression residual $g_2=f_2-cX_2$, linearity
gives the exact cohort decomposition $g_2=\sum_S g_S$.
Equivalently the coefficient in the primary linearization is
\begin{equation}
 A_+-c=-(c_\Sigma+c_{\KK}+c_{\BB}+\tau_+).
 \label{an03cfrv1:ledgerdifference}
\end{equation}
Thus global orthogonality is never used to remove a selector.
Combining the proved estimates, whenever the stated core
geometry holds at the state in question, gives
\begin{equation}
 \left\|\E_{P_2}[g_2X_2I_x]\right\|_{1,\infty}
 \le C_{L_c,R}\left[
 q\sqrt{E_{\Gp}J_{\Gp}}+\mu_{\Gp}\mathcal T_+
 +q^2\mu_{\KK}
 +E_\Sigma+q\sqrt{E_\Sigma J_\Sigma}+\mu_{\BB}\right].
 \label{an03cfrv1:totaldefect}
\end{equation}
This proves the selected-moment part of the suggested
decomposition with its unpaid secondary term explicit.

\section{Non-core forcing and its sign}
Let $O=\Gp\cup\Sigma\cup\BB$. Natural receiver-independent
original-state moments for CLOCK's weak terms are
\begin{equation}
 V_O=-\frac{\E_{P_2}[(f_O)_1X_2]}{qQ},\qquad
 K_O=\frac{\E_{P_1}[(f_O)_2X_1]}q .
 \label{an03cfrv1:VK}
\end{equation}
These are definitions, not a claim that all remaining
non-core contributions have already been bounded.
Indeed write $(f_O)_1=-qV_OX_2+g_{O,1}$ with
$\E_{P_2}[g_{O,1}X_2]=0$. Its linear part produces
$\lambda V_O\Phi(\rho x)$ in $2f_x$ to leading order;
the selected moments of $g_{O,1}$ must still be charged.
The full $P_1$ moment produces $-K_O$ in $2f_y$.
This fixes the signs and the mixture factors in the proposed
identification. It agrees with LOC's moments for weak-chart
donors $U/\sqrt m=(qu,1)+O(q^2)$,
$W/\sqrt m=(-qv,1)+O(q^2)$.

\begin{lemma}[Negative part and the uncharged secondary size]
\label{an03cfrv1:negative}
Under the primary signs from RCET,
\begin{align}
 (K_O)_-&\le C\left[
 q^{-1}\mu_{\Gp}e^{-c/q^2}
 +E_\Sigma+\frac{\sqrt{E_\Sigma J_\Sigma}}q
 +\frac{\mu_{\BB}}q\right],\label{an03cfrv1:negativebound}\\
 |V_\Sigma|&\le C\left[
 J_\Sigma+\frac{\sqrt{E_\Sigma J_\Sigma}}q\right].
 \label{an03cfrv1:Vsecondary}
\end{align}
The constants depend only on the ratio box and Gaussian moments.
One may define $K_w=(K_O)_+\ge0$, but this transfers
$(K_O)_-/2$ into the $y$-remainder rather than eliminating it.
\end{lemma}
\begin{proof}
Since $(f_{\Gp})_2\ge0$, its $P_1$ moment is negative only
where $X_1<0$. On that set
$|(f_{\Gp})_2X_1|\le\mu_{\Gp}|X||X_1|$.
A Gaussian tail bound with mean $1$ gives its first term in
\eqref{an03cfrv1:negativebound}, after division by $q$.
For any cohort, Minkowski and balance give
\begin{align*}
 \|(f_S)_2\|_{P_1}
 &\le2\sqrt{E_SJ_S}+\sqrt2q\sqrt{E_S(E_S+J_S)},\\
 \|(f_S)_1\|_{P_2}
 &\le2\rho\sqrt{E_SJ_S}+\sqrt2q\sqrt{J_S(E_S+J_S)}.
\end{align*}
Cauchy--Schwarz with $X_1$ or $X_2$, respectively, proves
the secondary bounds, using $q\le1$ and
$\sqrt{E(E+J)}\le E+\sqrt{EJ}$.
For $\BB$, $|f_{\BB}|\le\mu_{\BB}|X|$ gives cost
$C\mu_{\BB}/q$. Finally
$-K_O=-(K_O)_++(K_O)_-$ proves the remainder assertion.
\end{proof}

Under the stated guards the secondary quantities exposed in
\eqref{an03cfrv1:totaldefect} and
\eqref{an03cfrv1:negativebound} are bounded above by a constant
times $E_\Sigma+K_Y\sqrt{E_\Sigma}$, apart from vanishing
$q^2K_Y^2$ terms. There is no claim that these bounds are
sharp lower bounds on the actual field.

\section{A radial state identity, with the missing budget visible}
\begin{proposition}[Radial reduction at a charted core state]
\label{an03cfrv1:radial}
At any state satisfying the residual guards, suppose the core
donors have $|\theta|,|\psi|\le Rq$ and
$\mu_{\KK}\le\overline M_{\KK}$, with fixed $R,\overline M_{\KK}$.
For a receiver with $|\theta|,|\psi|\le L_cq$, set
$\kappa_1=\int_{\KK}W_1U_1$.
Then
\begin{equation}
 (\log m)'=1-\kappa_1+\delta_m,\qquad
 |\delta_m|\le C_{L_c,R}\left[
 q^2(1+\overline M_{\KK}+M_{\rm res})
 +J_G+q\sqrt{E_GJ_G}+\mu_{\BB}\right].
 \label{an03cfrv1:radialbound}
\end{equation}
The constant also uses the guard-independent diagonal loss bound.
This is a state lemma; continued core geometry and a uniform
mass budget are not inferred from the four residual guards.
\end{proposition}
\begin{proof}
On $P_1$, a core donor satisfies
\[
 \|(f_{\KK})_1-\kappa_1X_1\|_{P_1}
 \le C_Rq^2\mu_{\KK}.
\]
To check this, replace its positive part by $U\cdot X$,
retaining a Gaussian negative-gate remainder. The retained
$X_2$ coefficient has size at most $C_Rq\mu_{\KK}$ and
$\|X_2\|_{P_1}=q$; the gate remainder is exponentially small.
For $G$, Minkowski gives
$\|(f_G)_1\|_{P_1}\le C[J_G+q\sqrt{E_GJ_G}]$.
The outside norm is at most $C\mu_{\BB}$.

Removing the \emph{receiver} gate from the $P_1$ residual
moment costs $C_{L_c}e^{-c/q^2}$: use the diagonal loss bound
and the second moment of $X_1$ on the receiver's negative
half-space. The $P_2$ residual contribution to $R_{11}$
is at most
$\|e_1\|_{P_2}\|X_1\|_{P_2}/2\le M_{\rm res}q^2/2$.
It follows that
\[
 2R_{11}=1-\kappa_1+
 O_{L_c,R}\big(q^2(1+\overline M_{\KK}+M_{\rm res})
       +J_G+q\sqrt{E_GJ_G}+\mu_{\BB}\big).
\]
For unit directions $u,w$, the exact identity is
$(\log m)'=2w^TR_\theta u$.
The chart gives $w_1u_1=1+O(L_c^2q^2)$ and
$|u_2|,|w_2|\le L_cq$.
BULK's residual-entry estimate bounds cross entries by
$C(1+M_{\rm res})q$; the diagonal entries are bounded by the
loss constant plus $CM_{\rm res}q^2$.
Expansion proves \eqref{an03cfrv1:radialbound}.
\end{proof}

If this state geometry holds on the stopped interval, the
displayed bound integrates to
\begin{equation}
 \int|\delta_m|\le C_{L_c,R}\left[
 q^2(1+\overline M_{\KK}+M_{\rm res}+K_Y^2)T_*
 +q^2K_Y\int\sqrt{E_g+E_\Sigma}
 +\int\mu_{\BB}\right].
 \label{an03cfrv1:radialintegral}
\end{equation}
This identifies the radial input that a later $B_s$ proof
would have to close. It neither proves $B_s$ nor establishes
a strong mass law on this interval.

\section{Exact missing inequalities and stop}
The suggested secondary charge would be justified by the
following all-endpoint inequality:
\begin{equation}
 \boxed{\quad E_\Sigma(t)\le
 C_\Sigma\{q^5+q^{\Gamma_\nu}H(t)\}
 \qquad(t_0\le t\le t_*).\quad}
 \label{an03cfrv1:missingenergy}
\end{equation}
What SEP actually gives here is the entry profile and
$E_\Sigma(t_0)=O(q^5)$, not
\eqref{an03cfrv1:missingenergy}.
The source proof of the latter explicitly invokes both
\begin{align}
 D_\beta(t)&\le C_D\xi_\beta e^{\beta_0I(t)/2}
       \quad\hbox{until }qD_\beta=\delta_0,\label{an03cfrv1:missingdistance}\\
 m(t,\beta)&\le C_m m(t_0,\beta)
       \max\{1,C_gq^2\xi_\beta^2e^{\beta_0I(t)}\}^{1+\nu}
       \quad\hbox{at every endpoint}.\label{an03cfrv1:missinggain}
\end{align}
The symbol $\beta_0$ denotes the gain/distance exponent, not
the initial label. Neither display is a consequence already
proved in the cited sources under \eqref{an03cfrv1:guards}.
They are the open S2 contracts identified by roadmap v1.8.

For the estimates proved above, the missing integrated
secondary budget is explicitly
\begin{equation}
 \boxed{\quad
 \mathcal U_\Sigma(t_*):=
 \int_{t_0}^{t_*}
       (E_\Sigma+K_Y\sqrt{E_\Sigma})\dd t=o_q(1).
 \quad}
 \label{an03cfrv1:missingintegral}
\end{equation}
It would follow immediately from
\eqref{an03cfrv1:missingenergy}, since $H\le E_{\Gp}<E_g$:
\[
 \mathcal U_\Sigma\le
 C T_*\left[q^5+q^{\Gamma_\nu}E_g
       +K_Y(q^{5/2}+q^{\Gamma_\nu/2}\sqrt{E_g})\right]=o(1).
\]
This calculation explains the useful positive powers of $q$,
but is not a proof that its premise holds.
The actual guards control $E_{\Gp}$ and $J_{\Gp}+J_\Sigma$,
not $E_\Sigma$. Loss controls the total output residual,
not a cancellation-free cohort energy. No assertion that
arbitrary energy functions are realized by initialized
trajectories is needed to locate this source dependency.
A different proof that bypasses the secondary budget is not
ruled out; none is supplied by the proposed source implication.

There is also a separate outside-population bookkeeping gap
for the requested complete value/derivative remainder.
The generic PROFILE/MOMC charge is $C\mu_{\BB}/q$, so this
route needs
\begin{equation}
 \int_{t_0}^{t_*}\mu_{\BB}/q=o_q(1)
 \quad\hbox{or an explicitly allowed small-fixed budget}.
 \label{an03cfrv1:missingoutside}
\end{equation}
LCRC retains only $\sup\mu_{\BB}/q=o(1)$, which on an
$O(\ell)$ interval does not imply this display.
For example the scalar budget $\mu_{\BB}/q=\ell^{-1/2}$
on an interval of length $\ell$ has vanishing supremum
and integral $\sqrt\ell$.
This is a counterexample to a norm-budget inference,
\emph{not} an asserted network trajectory.
MOMC's stronger integrated outside bound stops at $t_h$;
it cannot be reused after $t_0=t_h$.
The core geometry and mass bounds in LCRC are also retained
state inputs, not conclusions of the four residual guards.
Even granting them does not repair
\eqref{an03cfrv1:missingenergy}.

Finally LCRC's sharper cap $H_{\Gp}\le(9/8)c_b+o(1)$ also
uses its secondary estimate in the ledger cap. It cannot
be imported here solely from the energy guard to conclude
$\int c=O(c_b)$.
Thus we have not established a full c-form field with
nonnegative $K_w$ and
\[
 \int\epsilon_F\le\kappa(C_0,c_b)+o_q(1),\qquad
 \kappa(C_0,c_b)\longrightarrow0 .
\]
The closed primary bound \eqref{an03cfrv1:primaryintegral}
is a subset of that target, not a substitute for it.
Per the stop rule, we stop at the unlisted S2 dependency.
We do not assume \eqref{an03cfrv1:missingdistance} or
\eqref{an03cfrv1:missinggain}, evolve the complement moment,
apply CLOCK's passage theorem, or start selection.

\section{Constant audit}
The fixed receiver radius $L_c$ and donor radius $R$ are
chosen before residual guards.
Constants $C_{L_c,R}$ use only these radii, the compact
$\rho$-box, the initial-loss bound, and Gaussian moments.
A core mass budget $\overline M_{\KK}$, when explicitly
used in the radial state lemma, must be supplied independently
of the residual guards; this document does not construct it.
There is no replacement of these constants by a radius
depending on $M_{\rm res}$.

RCET's $C_*,C_{\rm RCET}$ and the admissible smallness of
$C_0$ may depend on $(B,M_{\rm res},c_b)$.
In the closed estimate they multiply $C_0$ or vanishing
$q$-scales. $E_g$ and $K_Y$ likewise occur only as coefficients
of the displayed small scales in the primary estimate.
They do not furnish a small secondary or outside budget.
The quantities $V_O,K_O$ are state moments, not uniform
constants. Their uncontrolled pieces are shown explicitly,
rather than counted as guard-independent order-one constants.

\begin{thebibliography}{9}\small
\bibitem{Roadmap} Roadmap version 1.8, S0 assembly requirements,
S2 and S3 (Oct.\ 21 refined route).
\bibitem{CLOCK} \path{AN03_exact_amplitude_clock_and_c_passage_v1.tex},
\texttt{inc:AN03:clock:v1:cfield}, \texttt{passage}, Section 6.
\bibitem{LCRC} \path{AN03_ledger_cap_residual_closure_v1.tex},
Section 1; \texttt{an03lcrcv1:secondary},
\texttt{secondarybound}, \texttt{restsmall}, \texttt{closure}.
\bibitem{RCET} \path{AN03_residual_closure_and_eta0_transport_v1.tex},
Theorem 2.1 and its growing-margin proof; the conservative
$y$-comparison is used as already corrected in LCRC.
\bibitem{BULK} \path{AN03_bulk_clock_and_secondary_Q2_energy_v1.tex},
Theorem 4.1, Proposition 5.2, secondary corollary,
residual-entry bounds.
\bibitem{MOMC} \path{AN02_moving_overlap_M_clock_passage_v1.tex},
finite-noise reduction and its Gaussian boundary proof;
outside budgets only through $t_h$.
\bibitem{Entry} \path{AN02_Q2_entry_contracts_v1.tex} and
\path{AN02_secondary_entry_profile_v1.tex};
primary entry and secondary entry profile, respectively.
\bibitem{LOC} \path{AN03_local_rate_mean_tracking_v1.tex},
leading field and moments, Lemma 3.1; no original-flow
future budget is inferred from the leading passage.
\end{thebibliography}
\end{document}

